\chapter{Аналитическая часть}

\section{Постановка задачи конвейерной обработки заявок}

В лабораторной работе требуется реализовать конвейер обработки заявок, где каждая заявка соответствует одному входному файлу в формате Graphviz (dot).
Заявка на входе содержит имя файла, а в процессе обработки пополняется данными и результатами вычислений.

Конвейер состоит из трёх обслуживающих устройств (ОУ), работающих в отдельных нативных потоках, и взаимодействующих через очереди:
\begin{itemize}
	\item оу1 --- чтение входных данных из файла и заполнение структуры заявки;
	\item оу2 --- выполнение вычислительного задания по варианту ЛР5, при этом внутри ОУ2 используется оптимальное число вспомогательных потоков $k$, определённое в ЛР5;
	\item оу3 --- дамп результата в выходной файл (одна заявка --- один файл), имя которого получается добавлением постфикса к исходному имени.
\end{itemize}

Перед запуском конвейера входная очередь, предназначенная для ОУ1, должна быть заполнена $N$ заявками, где
\[
N \in \{25, 50, 75, 100, 125\}.
\]
Требуется сравнить реальное время обработки этих $N$ заявок в двух режимах:
\begin{itemize}
	\item последовательная обработка;
	\item конвейерная обработка на трёх нативных потоках (или го-рутинах), не считая диспетчера.
\end{itemize}

\section{Модель заявки}

Заявка представляет собой объект, который проходит последовательно через этапы конвейера и постепенно «обогащается» данными.
Для формального описания можно представить заявку как кортеж
\begin{equation}
	req = \langle name,\, in,\, res,\, out \rangle,
	\label{eq:req_model}
\end{equation}
где:
\begin{itemize}
	\item $name$ --- имя входного файла (dot);
	\item $in$ --- данные, полученные на этапе чтения (например, представление графа в памяти);
	\item $res$ --- результат вычислений, полученный на этапе обработки по варианту ЛР5;
	\item $out$ --- данные, предназначенные для записи в выходной файл (например, текст/структуры, полученные из $res$).
\end{itemize}

Идея «пополнения заявки» означает, что каждый этап конвейера добавляет или модифицирует поля заявки, не разрушая уже полученные данные~\cite{valedinsky2018}.

\section{Обслуживающие устройства конвейера}

Работа конвейера описывается как последовательное применение преобразований к заявке:
\begin{equation}
	req_1 = F_1(req), \quad req_2 = F_2(req_1), \quad req_3 = F_3(req_2),
	\label{eq:ou_funcs}
\end{equation}
где $F_1, F_2, F_3$ соответствуют ОУ1, ОУ2 и ОУ3 соответственно.

\subsection{ОУ1: чтение входных данных}
ОУ1 выполняет чтение входного файла, указанного в поле $name$, и формирует представление входных данных в памяти.
Результат чтения записывается в поле $in$ заявки.
Данный этап включает операции ввода-вывода и преобразование данных файла в структуру, пригодную для дальнейшей обработки.

\subsection{ОУ2: вычислительная обработка по варианту ЛР5}
ОУ2 выполняет вычислительную часть, соответствующую варианту ЛР5, используя данные $in$.
Результат вычислений сохраняется в поле $res$.

Согласно заданию, внутри ОУ2 используется оптимальное число вспомогательных потоков $k$, определённое в ЛР5.
Таким образом, ОУ2 является этапом, который может включать внутренний параллелизм: один поток ОУ2 организует работу группы из $k$ рабочих потоков и собирает итоговый результат~\cite{cpp_standard}.

\subsection{ОУ3: запись результата}
ОУ3 формирует выходной файл для каждой заявки и выполняет запись результата.
Имя выходного файла получается добавлением постфикса к имени исходного файла, заданного в $name$.
Результаты записи относятся к одной заявке (один входной файл --- один выходной файл), что упрощает сопоставление результатов обработки.

\section{Взаимодействие между этапами через очереди}

Обслуживающие устройства обмениваются заявками через очереди.
Пусть $Q_{12}$ --- очередь между ОУ1 и ОУ2, а $Q_{23}$ --- очередь между ОУ2 и ОУ3.
Также существует входная очередь $Q_{in}$, из которой ОУ1 получает исходные заявки.

Общая схема взаимодействия:
\begin{itemize}
	\item до запуска конвейера в $Q_{in}$ помещаются $N$ заявок;
	\item оу1 извлекает заявку из $Q_{in}$, выполняет $F_1$ и помещает результат в $Q_{12}$;
	\item оу2 извлекает заявку из $Q_{12}$, выполняет $F_2$ (с использованием $k$ вспомогательных потоков) и помещает результат в $Q_{23}$;
	\item оу3 извлекает заявку из $Q_{23}$, выполняет $F_3$ и завершает обработку заявки.
\end{itemize}

Так как очереди являются разделяемыми структурами данных, для корректной работы конвейера требуется синхронизация доступа к ним.
В модели нативных потоков для этого применяются стандартные средства синхронизации: мьютексы, условные переменные и/или атомарные операции~\cite{cpp_standard}.

\section{Последовательная обработка заявок}

Последовательный режим служит базовой линией сравнения.
В этом режиме заявки обрабатываются одна за другой в одном потоке, причём для каждой заявки выполняются все этапы в фиксированном порядке:
\begin{enumerate}
	\item чтение данных (аналог ОУ1);
	\item вычисление результата (аналог ОУ2, при необходимости с использованием $k$ потоков как в варианте);
	\item запись результата (аналог ОУ3).
\end{enumerate}

Главная особенность последовательного режима состоит в отсутствии перекрытия этапов: чтение следующей заявки не начинается, пока полностью не завершены вычисления и запись результата по текущей заявке.

\section{Конвейерная обработка заявок}

Конвейерный режим реализует перекрытие этапов обработки во времени:
пока ОУ2 выполняет вычисления для одной заявки, ОУ1 может читать следующую, а ОУ3 --- записывать результат для предыдущей.
За счёт этого повышается общая пропускная способность системы при обработке серии заявок.

Эффективность конвейера зависит от согласованности скоростей этапов:
если один из этапов существенно медленнее остальных, он становится «узким местом», и очереди между этапами начинают накапливать заявки.

\section{Дополнительное обслуживающее устройство ОУ+ (ЛР8)}

В дополнительном задании ЛР8 между ОУ2 и ОУ3 добавляется обслуживающее устройство ОУ+, то есть конвейер принимает вид:
\[
\text{ОУ1} \rightarrow \text{ОУ2} \rightarrow \text{ОУ+} \rightarrow \text{ОУ3}.
\]

ОУ+ выполняет формирование dot-файла, в котором компоненты связности окрашены в цвета (одна компонента --- один цвет).
С точки зрения конвейера, ОУ+ является отдельным этапом, который:
\begin{itemize}
	\item получает из ОУ2 результаты анализа (например, разбиение на компоненты связности) в поле $res$;
	\item формирует представление для визуализации (обновлённое описание графа в формате dot) и передаёт его дальше на запись.
\end{itemize}

После добавления ОУ+ требуется ввести дополнительную очередь между ОУ2 и ОУ+, а также учитывать, что добавленный этап также может стать ограничивающим по скорости.

\section{Основные понятия многопоточности, используемые в работе}

При реализации конвейера на нативных потоках применяются следующие понятия~\cite{cpp_standard}:
\begin{itemize}
	\item поток выполнения --- независимая последовательность исполнения в рамках процесса;
	\item разделяемые данные --- области памяти (в том числе очереди), доступные нескольким потокам;
	\item критическая секция --- участок кода, в котором требуется взаимное исключение;
	\item мьютекс --- механизм взаимного исключения для защиты критических секций;
	\item условная переменная --- механизм ожидания события (например, появления элемента в очереди).
\end{itemize}

Использование данных механизмов обеспечивает корректную передачу заявок по конвейеру и предотвращает гонки данных при параллельном выполнении этапов.

\section*{Вывод}

В аналитической части сформулирована задача конвейерной обработки заявок, описаны структура заявки и назначение обслуживающих устройств ОУ1--ОУ3.
Рассмотрены два режима работы: последовательный и конвейерный с передачей данных через очереди и синхронизацией между нативными потоками.
Также описано дополнительное устройство ОУ+ из ЛР8, формирующее dot-файл с раскраской компонент связности.
Полученные положения служат основой для проектирования структуры программы и реализации конвейера~\cite{valedinsky2018,cpp_standard}.
